% flatex input: [main.tex]
\documentclass[oneside, 12pt]{book}
% \usepackage[slovak]{babel}
\usepackage{epsfig}
\usepackage{color}
\usepackage{url}
\usepackage[utf8]{inputenc}
\usepackage[T1]{fontenc}
\usepackage{lmodern}
\usepackage{todonotes}
\usepackage{amsfonts}
\reversemarginpar
\usepackage[]{algorithm2e}

% flatex input: [./config.tex]

%------------------------------------------------------------------------------------------------------
% pekne pokope definujeme potrebne udaje
%------------------------------------------------------------------------------------------------------
\def\mftitle{Two-way databinding of models and views in Dart}          % nazov diplomovky/bakalarky
\def\mfthesistype{Diploma thesis}             % typ prace
\def\mfkeywords{Databinding, Dart, Facebook React, AngularJS, ... }        % proste klucove slova idealne oddelene ciarkou
\def\mfkeywordsen{Databinding, Dart, Facebook React, AngularJS, ... }      % klucove slova v anglictine EN
\def\mfkeywordssk{Databinding, Dart, Facebook React, AngularJS, ... }      % klucove slova v slovencine SK
\def\mfauthor{Bc. Jakub Uhrík}                            % vase meno a priezvysko
\def\mfadvisor{RNDr. Tomáš Kulich, PhD.}                                  % cele meno skolitela
\def\mfplacedate{Bratislava, 2014}                        % mesto a rok
\def\mfdate{2014}                                         % Rok
\def\mfuniversity{Comenius University in Bratislava}      % Univerzita
\def\mffakulta{Faculty of Mathematics, Physics and Informatics}  % Katedra
\def\mffakultafull{Faculty of Mathematics, Physics and Informatics, Comenius University in Bratislava}  
                                                          % Katedra v plnom zneni aj s univerzitou
\def\mfprogram{Computer Science}                               % studijny program
\def\mfodbor{9.2.1. Computer Science, Informatics}                            % studijny odbor
\def\mfpracovisko{FMFI.KI - Department of Computer Science}                    % skoliace pracovisko


% flatex input end: [./config.tex]

% \usepackage[slovak]{babel}

\usepackage[a4paper,vmargin={25mm,25mm},hmargin={35mm,20mm}]{geometry} 
\usepackage{minted}
\usepackage{caption}
\hyphenpenalty=10000
\tolerance=10000

\usepackage[
	bookmarks=true,
	bookmarksnumbered=false, 	% true means bookmarks in
								% left window are numbered
	bookmarksopen=false,
								% true means only level 1
								% are displayed.
	colorlinks=true,
	pdftitle={\mftitle},
	pdfauthor={\mfauthor},
	pdfsubject={\mfthesistype},
	pdfkeywords={\mfkeywords},
	pdfproducer={\mffakultafull},
	linkcolor=webblue]{hyperref}
\definecolor{webgreen}{rgb}{0, 0.5, 0} % less intense green
\definecolor{webblue}{rgb}{0, 0, 0.5} % less intense blue
\definecolor{webred}{rgb}{0.5, 0, 0}
\definecolor{grey}{rgb}{0.5, 0.5, 0.5}
\definecolor{white}{rgb}{0.95, 0.95, 0.95}
\definecolor{black}{rgb}{0, 0, 0}
%------------------------------------------------------------------------------------------------------
% less intense red
%------------------------------------------------------------------------------------------------------

\usepackage[final]{pdfpages}

 
%------------------------------------------------------------------------------------------------------
% riadkovanie jeden a pol
%------------------------------------------------------------------------------------------------------
\renewcommand\baselinestretch{1.3} 

\def\phrase{\begingroup\phrasenohyphen
	\dononbreakablespace
	\dononbreakablehyphen
	\dophrase}
\makeatletter
\def\phrasenohyphen{%
	\language\@ifundefined{l@nohyphenation}{\@cclv}{\l@nohyphenation}}
\makeatother
\def\dophrase#1{#1\endgroup}

\newcommand*{\fullref}[1]{\hyperref[{#1}]{\autoref*{#1} \nameref*{#1}}}

\author{\mfauthor}
\title{\mftitle}

\begin{document}

%------------------------------------------------------------------------------------------------------
% vlozime strukturu, vyplnte ju priamov suboroch, treba len pomenit texty
%------------------------------------------------------------------------------------------------------
% flatex input: [./content/structure/obalka.tex]
%------------------------------------------------------------------------------------------------------
%  toto je subor obalky vasej prace. 
%  s tymto by ste asi ani nemali nic robit pouziju sa prednastavene veci z hlavneho suboru
%------------------------------------------------------------------------------------------------------
\frontmatter

\thispagestyle{empty}

\noindent
\begin{center}
\begin{minipage}{0.8\textwidth}
\centerline{\renewcommand\baselinestretch{1.3} \LARGE\sc\mfuniversity}
\centerline{\sc\mffakulta}
\end{minipage}
\end{center}

\vfill
\begin{center}
\begin{minipage}{1\textwidth}
\bigskip\bigskip
\begin{center}
\linespread{1}\LARGE\sc\mftitle
\end{center}
\smallskip
\centerline{\mfthesistype}
\bigskip
\bigskip
\bigskip\bigskip
\end{minipage}
\end{center}
\vfill
{\bf\mfdate\\
\indent\mfauthor}
\eject 


% flatex input end: [./content/structure/obalka.tex]

%------------------------------------------------------------------------------------------------------

% flatex input: [./content/structure/titulka.tex]
%------------------------------------------------------------------------------------------------------
%  toto je titulna stranka, vsetky veci ste nastavili v config.tex
%------------------------------------------------------------------------------------------------------
\thispagestyle{empty}

\noindent
\begin{center}
\begin{minipage}{0.8\textwidth}
\centerline{\LARGE\sc\mfuniversity}
\centerline{\sc\mffakulta}
\end{minipage}
\end{center}

\vfill
\begin{center}
\begin{minipage}{1\textwidth}
\bigskip\bigskip
\begin{center}
\linespread{1}\LARGE\sc\mftitle
\end{center}
\smallskip
\centerline{\mfthesistype}
\bigskip
\bigskip
\bigskip\bigskip
\end{minipage}
\end{center}
\vfill
\begin{minipage}{0.8\textwidth}
\begin{tabular}{l l}
Study programme:& \mfprogram \\
Field of Study:& \mfodbor \\
Department:& \mfpracovisko\\
Thesis supervisor:&   \mfadvisor \\
\end{tabular}
\end{minipage}
\begin{center}
\end{center}
\vfill
{\bf\mfplacedate\\
\indent\mfauthor}
\eject 


% flatex input end: [./content/structure/titulka.tex]

%------------------------------------------------------------------------------------------------------

% flatex input: [./content/structure/zadanie.tex]
%------------------------------------------------------------------------------------------------------
%  pdf zadania dajte do zlozky images v subore zadanie.pdf tak ako vidite v prikaze. 
%  ak zadavate aj anglicke zadanie, odkomentujte(odstrante percento zo zaciatku) 
%  druhy riadok a podla neho pridajte adekvatne pdf
%------------------------------------------------------------------------------------------------------

\includepdf{images/zadanie-en.pdf}
\includepdf{images/zadanie.pdf}
% \includepdf{images/zadanie-en.pdf}


% flatex input end: [./content/structure/zadanie.tex]

%------------------------------------------------------------------------------------------------------

%------------------------------------------------------------------------------------------------------
% ked treba, da sa vlozit cestne prehlasenie...
%------------------------------------------------------------------------------------------------------
%%------------------------------------------------------------------------------------------------------
% prehlasenie, ze ste to poctivo sami vypracovali...
%------------------------------------------------------------------------------------------------------
{~}\vspace{12cm}

\noindent
\begin{minipage}{0.25\textwidth}~\end{minipage}
\begin{minipage}{0.68\textwidth}
%------------------------------------------------------------------------------------------------------
% sem kludne pridajte vase originalne cestne vyhlasenie, ale mozete sa inspirovat mojim
%------------------------------------------------------------------------------------------------------
I hereby declare that this thesis represents my own work and effort. Where other sourcesof information have been used, they have been acknowledged.

\bigskip\bigskip

\hfill\hbox to 6cm{\dotfill}
\end{minipage}
\vfill\eject 



% flatex input: [./content/structure/thanks.tex]
%------------------------------------------------------------------------------------------------------
% Podakovanie. Proste subor s podakovanim. Napiste sem komu dakujete, mozete sa inspirovat mojim
%------------------------------------------------------------------------------------------------------
{~}\vspace{12cm}

\noindent
\begin{minipage}{0.25\textwidth}~\end{minipage}
\begin{center}
\begin{minipage}{1\textwidth}
%------------------------------------------------------------------------------------------------------
%  sem pridajte svoj originalny text, alebo len prepiste \mfadvisor za vysklonovane meno vaseho veduceho
%------------------------------------------------------------------------------------------------------

I would like to thank my supervisor \mfadvisor for his guidance,support, and encouragement throughout writing this thesis. 

Special thanks belong to my family for all their support.

\end{minipage}
\end{center}
\hfill\mfauthor
\vfill\eject 


% flatex input end: [./content/structure/thanks.tex]

%%------------------------------------------------------------------------------------------------------
% prehlasenie, ze ste to poctivo sami vypracovali...
%------------------------------------------------------------------------------------------------------
{~}\vspace{12cm}

\noindent
\begin{minipage}{0.25\textwidth}~\end{minipage}
\begin{minipage}{0.68\textwidth}
%------------------------------------------------------------------------------------------------------
% sem kludne pridajte vase originalne cestne vyhlasenie, ale mozete sa inspirovat mojim
%------------------------------------------------------------------------------------------------------
I hereby declare that this thesis represents my own work and effort. Where other sourcesof information have been used, they have been acknowledged.

\bigskip\bigskip

\hfill\hbox to 6cm{\dotfill}
\end{minipage}
\vfill\eject 



% flatex input: [./content/structure/abstract.tex]
%------------------------------------------------------------------------------------------------------
% Toto je subor s abstraktom. Pridajte sem vas/tvoj abstrakt v slovencine a/alebo anglictine
% Ak chcete len jeden abstrakt, vymazte vsetko medzi komentarmi "jazyk zaciatok" a "jazy koniec"
%------------------------------------------------------------------------------------------------------

%------------------------------------------------------------------------------------------------------
% EN zaciatok
%------------------------------------------------------------------------------------------------------
\noindent
\begin{center}
\begin{minipage}{1\textwidth}
\centerline{\large Abstract}
Abstract in english.
%------------------------------------------------------------------------------------------------------
% sem pridajte abstrakt v anglictine
%------------------------------------------------------------------------------------------------------
\\ \\ 
{\bf Key words:} \mfkeywordsen
\end{minipage}
\end{center}
\eject % EOP v
%------------------------------------------------------------------------------------------------------
% EN konie
%------------------------------------------------------------------------------------------------------

%------------------------------------------------------------------------------------------------------
% SK zaciatok
%------------------------------------------------------------------------------------------------------
\noindent
\begin{center}
\begin{minipage}{1\textwidth}
\centerline{\large Abstrakt}
%------------------------------------------------------------------------------------------------------
% sem pridajte abstrakt v slovencine
%------------------------------------------------------------------------------------------------------
Abstrakt v slovencine.
\\ \\ 
{\bf Kľúčové slová:} \mfkeywordssk
\end{minipage}
\end{center}
\eject % EOP v
%------------------------------------------------------------------------------------------------------
% SK konie
%------------------------------------------------------------------------------------------------------


% flatex input end: [./content/structure/abstract.tex]

%%------------------------------------------------------------------------------------------------------
% prehlasenie, ze ste to poctivo sami vypracovali...
%------------------------------------------------------------------------------------------------------
{~}\vspace{12cm}

\noindent
\begin{minipage}{0.25\textwidth}~\end{minipage}
\begin{minipage}{0.68\textwidth}
%------------------------------------------------------------------------------------------------------
% sem kludne pridajte vase originalne cestne vyhlasenie, ale mozete sa inspirovat mojim
%------------------------------------------------------------------------------------------------------
I hereby declare that this thesis represents my own work and effort. Where other sourcesof information have been used, they have been acknowledged.

\bigskip\bigskip

\hfill\hbox to 6cm{\dotfill}
\end{minipage}
\vfill\eject 



%------------------------------------------------------------------------------------------------------
% obsah
%------------------------------------------------------------------------------------------------------
\tableofcontents

%------------------------------------------------------------------------------------------------------
% zoznam obrazkov
%------------------------------------------------------------------------------------------------------
\listoffigures

\mainmatter

%------------------------------------------------------------------------------------------------------
% page stype plain pre ciste strany bez cisiel stran pre uvod
%------------------------------------------------------------------------------------------------------
\pagestyle{plain}

%------------------------------------------------------------------------------------------------------
% a teraz includneme vsetky kapitoly, realne je to len subor, v ktorom sa incuduju kapitoly postupne
%------------------------------------------------------------------------------------------------------
% flatex input: [./content/macros.tex]
\newcommand{\tiles}[0]{\textbf{\texttt{Tiles}} }
\newcommand{\react}[0]{\textit{React} }
\newcommand{\facebook}[0]{Facebook }

% flatex input end: [./content/macros.tex]

%------------------------------------------------------------------------------------------------------
% flatex input: [./content/chapters/chapters.tex]
\pagestyle{plain}
% flatex input: [content/chapters/intro.tex]
\chapter*{Introduction}\label{chap:intro}
\addcontentsline{toc}{chapter}{Introduction}

As one of the results of this magister thesis is our new databinding library in dart, 
which is called \tiles. 
In next  text, we will use only \tiles to mention \textit{our new databinding library in dart}.

% flatex input end: [content/chapters/intro.tex]

\pagestyle{headings}
% flatex input: [content/chapters/motivation.tex]
\chapter{Motivation}\label{chap:motivation}

The first question, as always should be, is the motivation of this work. 
What is the motivation to create another library, that will handle databinding in Dart?

The motivation to create the \textbf{tiles} library contains from several aspects, 
which are not contain in no other \texttt{Dart} library
\footnote{We didn't find any suitable library and didn't hear about it}.

\begin{description}
	\item[Dart as a programming language] \hfill \\
		Dart language is young programming language with an active development and progress. 
		One of its advantages is optional typing, 
		the build-in compilation to the JavaScript, which enable programming a browser applications,
		and a Java-like virtual machine, which runs the Dart in the most commonly used operating systems.

		It is designed for the web applications with all necessary support for them. 
		As it enable the compilation into the JavaScript and running directly under the OS, 
		it also enable to share a source code between server part of an application 
		and its client application running in the web browser.

		Dart also guarantees browser compatibility, what is important for ease of web application development.
	\item[Testability] \hfill \\
		Very important aspect in a building complex application is the testability of the source code.

		Because of this, it is essential to use libraries, which enables easy testing and mocking components.
	\item[Server side rendering] \hfill \\
		Server side rendering is very important for user experience and for search engine optimization. 

		When we have the CPM\footnote{Computer Programming Language} which can be used as on the server, so in the clients browser, 
		it is natural to think about a use of the same source code to create an in-browser application, 
		and to render its page on the server. 
	\item[No templates] \hfill \\
		This aspect is important from two point of views: the testability and the server side rendering.

		From the testability point of view, it is easier to test and mock structures created in only one CPM. 
		If the template is used to create a component of the UI, 
		it is much more difficult to test it and also think about this testing.
		If this component is only one class in the CPM without dependences on another type of the information, 
		testing is more natural and easier to think about.

		From the server side rendering point of view, if we want to work with templates, 
		we need to access them differently when we work in the OS and in the clients browser.
		If we have the structure fully composed in one CPM, it is easier to compose the same HTML structure on the server as in the browser, 
		then if we have the structure composed by the template and the CPM.
	\item[Only one language] \hfill \\
		This aspect is very related with the previous one. 
		When the application is created fully in one programming language, 
		it is easier for programmers to work with it (they don't have to switch between different CPM).

		Also it is easier to compile whole application into the JavaScript, analyze the source code or refactor it.
	\item[Reliability] \hfill \\
		The reliability has significant importance in complex applications. 
		This reliability can be achieved by automatic tests, a robust design and quality development. 
\end{description}

When we take into account these aspects, there exists no library, which fulfill all of them.

As there is a need for this kind of a library, we decided to design and create one. 

% flatex input end: [content/chapters/motivation.tex]

% flatex input: [content/chapters/databinding.tex]
\chapter{Databinding}\label{chap:databinding}

In this chapter we will introduce the area of databinding more deeply then in introduction.

\section{One way databinding}\label{sec:onewaydatabinding}

Discuss one way databinding.

\section{Two way databinding}

Discuss two way databinding.

% flatex input end: [content/chapters/databinding.tex]

% flatex input: [content/chapters/existing.tex]
\chapter{Existing solutions}\label{chap:existing}

When we think about the building of the user interfaces, we can think about the building them from components.
The component is a part of the UI, which has a functionality, own look and maybe some interaction.

The HTML is a basic component structure. Every element is a component, all elements are composed into tree structure. 
Elements have some functionality, own look(e.g. image) and some of them have interactions(e.g. input).

So when a library want to bind the data with a view, it basically bind the data to a component.

There are different approaches of the creating these components and the connection between them and a data. 
Components can be created directly by a programming language (Component driven), 
or by using a template engine, which create components based on the template, 
which describe component structure, in the most cases by HTML-like syntax(Template driven).

These two approaches do the same thing, create structure of the components, different way.

\section{Template driven}\label{sec:existing-template}

  Template driven approach is, as the name predicts, based on the usage of a template engine. 
  Template engines take a template and the data and create a component structure, 
  which is reflected into the HTML representation of passed data in the form of the template. 
  They can be considered as a function $t:\mathbb D\mapsto\mathbb H$, 
  where $\mathbb D$ is a set of all possible data and $\mathbb H$ is a context-free language of valid HTML.

  An easy example of a template, for example using \textit{handlebars.js} can look like this one (from \textit{handlebars.js} website):
  % flatex input: [content/examples/handlebars_entry.tex]
\begin{minted}{html+django}
<div class="entry">
  <h1>{{title}}</h1>
  <div class="body">
    {{body}}
  </div>
</div>
\end{minted}

% flatex input end: [content/examples/handlebars_entry.tex]


  When programmer want to use this template, he should create data object, which minimal version in JSON format is in next example:
  % flatex input: [content/examples/handlebars_entry_data.tex]
\begin{minted}{json}
{
  "title": "Some title",
  "body": "This is the content of the page"
}
\end{minted}

% flatex input end: [content/examples/handlebars_entry_data.tex]


  When template is filled by this data, following HTML will be produced
  % flatex input: [content/examples/handlebars_entry_result.tex]
\begin{minted}{html}
<div class="entry">
  <h1>Some title</h1>
  <div class="body">
    This is the content of the page
  </div>
</div>
\end{minted}

% flatex input end: [content/examples/handlebars_entry_result.tex]


  Most of template engines also offer logic markup, which add possibility of the better control of a composed structure. 
  This is highly usable when programmer want to create more complex structures based on the data. 
  The typical example of this structure is the \texttt{<ul>} list generated from the array of items to render.
  
  This "in template" logic has on one hand some advantages, on the other hand, 
  the HTML syntax was not created to represent a logic, but an information.
  Because of this, more complex templates witch not so trivial logic in it becomes hard to read and understand.
  
  Easy use of the logic in the template is shown on the next example:
  % flatex input: [content/examples/handlebars_logic.tex]
\begin{minted}{html+django}
<h1>Comments</h1>

<div id="comments">
  {{#each comments}}
  <div class="entry">
    {{#if author}}
      <h1>{{firstName}} {{lastName}}</h1>
    {{else}}
      <h1>Unknown Author</h1>
    {{/if}}
    <div>{{body}}</div>
  </div>
  {{/each}}
</div>

\end{minted}

% flatex input end: [content/examples/handlebars_logic.tex]


  The template driven view are highly used because of the syntax similarity between the template and the resulting HTML. 
  It is easy to convert the HTML produced by a graphic designer to the template used in the source code. 
  Also programmers used to work in the HTML more easily write templates then some other representation of the component structure.

  Different libraries work with templates in a different way. 
  Some of them really parse the input template as a string, recognize component tree in it and work with the template that way.
  Others uses in-browser HTML parser to parse the template and then fill it with the data.
  This approach, because of its usage of tools accessible only from the browser, is more difficult to render on the server.

  Templates are mostly used two different ways:
  \begin{description}
    \item[Template used as View] \hfill \\
      Template is used to render HTML structure into some element. 
      Functionality of the HTML structure is then realized separately and attached to it.  
      This is used for example in the CanJS.
    \item[Template used as Component] \hfill \\
      The other (and more modern) use of template is to represent one component with attached functionality, 
      which can be represented later as custom HTML element in other templates. 
      In this template, other custom components can be created by using their custom HTML element representation. 
      It is not necessary to create them separately in the code of an application.

      This approach is used e.g. in Polymer project, which work with so called "shadow DOM" which use similar concept.
  \end{description}

  \fullref{table:template-driven-libraries} compares some of existing solutions which are standalone libraries or MVC frameworks.
  The aspects of the comparison are a natural rendering in the browser and  on the server and if the library is a standalone UI library, 
  or is a part of the more complex MVC framework. 
  We don't compare a possibility to render the view on the server other then the natural way, 
  because it is always possible to render it by usage of tools like the \textit{PhantomJS}.

  \begin{table}
    \begin{tabular}{|l|l|c|c|c|c|}
      \hline
      \textbf{Solution}& \textbf{Language}   & \textbf{Standalone} & \textbf{In browser} & \textbf{On server} \\
      \hline
      handlebars       & JavaScript          &         yes         &        yes          &        yes         \\
      \hline
      {{mustache}}     & \shortstack{JavaScript,  
                            \\    python...} &        yes         &        yes          &        yes          \\
      \hline
      dust             & JavaScript          &         yes         &        yes          &        yes         \\
      \hline
      AngularJS        & JavaScript          &         no          &        yes          &        no          \\
      \hline
      meteor           & JavaScript          &         no          &        yes          &        no          \\
      \hline
      EmberJS          & JavaScript          &         no          &        yes          &        yes         \\
      \hline
      Derby            & JavaScript          &         no          &        yes          &        yes         \\
      \hline
      Polymer          & JavaScript          &         yes         &        yes          &        no          \\
      \hline
      Polymer.dart     & Dart                &         yes         &        yes          &      not now       \\
      \hline
    \end{tabular}
    \caption{Comparison of template driven libraries}
    \label{table:template-driven-libraries}
  \end{table}

\section{Component driven}\label{sec:existing-component}

  Component driven views, in opposite to the template driven, don't use any additional type of data like templates. 
  Components are created by the same programming language as the functionality and 
  are composed into the tree structure which is mapped into the DOM. 

  When the tree of components (we will call it "Virtual DOM" later)is constructed, 
  it is rendered to the DOM by the depth-first search of the component tree.
  When components and HTML elements are connected by stored associations, 
  every change in the component structure can be applied to the DOM tree. 

  In addition, if we have the tree of components, we can easily, by the similar depth-first search, 
  create the markup string representing the HTML markup of the component tree.
  This enable the rendering of the whole component tree on the server without use of browser-specific features.

  An example of the component driven UI library is the JavaScript library \react created by the \facebook.
  \react is standalone UI library which enable native rendering of the component structure as in the browser, so on the server.

  We decided to use a similar approach to \react library, so we briefly describe it.

  \subsection{React}\label{subsec:existing-component-react}
    
    Lots of people use \react as the V in MVC.\cite{react}

    \react is JavaScript UI library from \facebook. 
    Its main concept is to pack parts of the web application into reusable components, 
    which are represented as object in JavaScript. 

    This components can be mounted into elements in DOM, for now, we will call it \textit{mount root}. 
    This will create \textit{virtual DOM} "mounted" to \textit{mount root}. 
    This virtual DOM is then reflexed into the real DOM under the \textit{mount root}.

    \react uses a virtual DOM diff implementation for ultra-high performance. 
    It can also render on the server using Node.js — no heavy browser DOM required.\cite{react}

    Components are organized to the virtual DOM tree, where a data flows from the root component to leaves. 
    This data flow is implemented by the props of the component, which are read-only.
    Component have also a state, which should be stored in the state attribute and updated by methods \texttt{setState} and \texttt{replaceState}. 
    The state shouldn't be updated directly to preserve the invariant, that the real DOM always represents the actual state of the virtual one. 

    Component describe the structure under it by its method \texttt{render}, 
    which should return one instance of a component, 
    which will be added as a child of this component.
    The render function also add props to the child component.
    This realize the data flow in "down" direction. 
    The render also add children to the child component, which is the way, 
    how to create a spreading tree, not just a line. 
    The child component have read access to passed children and can reuse them in the \texttt{render} method or ignore them.

    The \react offer own events system with synthetic event bubbling.
    This enable programmer to listen to events independently from browser. 
    The \react manage the browser compatibility.

    Components can listen to events on DOM components (internal \react components, representing DOM elements). 
    They are attached trough props by the event listeners syntax(\texttt{onChange}, \texttt{onClick} etc.). 

    State change(by mentioned methods) trigger redrawing of virtual DOM. 

    \react implements component life-cycle methods,
    which notify the component about its actual state of living (just mounted, just updated, before unmount, etc.). 
    They are the superset of life-cycle methods implemented in the tiles library. 

    For more information about \react, it's architecture and API 
    reader can go to the website of the \react project \url[http://facebook.github.io/react/].

\section{Conclusion}\label{sec:existing-conclusion}

  We decided to use \textbf{Component driven} views and databinding, because it is not dependent on the template engine, 
  whole source code can be written in the same CPL with all advantages gained by that 
  and naturally easier thinking about testing and mocking.

  Our solution is based on the idea of the \react library.

  As we decided to work in the Dart language, we don't have to implement a browser compatibility, synthetic events, mixins, etc.

  The architecture of the tiles library will be described in the \fullref{chap:oursolution}.

% flatex input end: [content/chapters/existing.tex]

% flatex input: [content/chapters/our.tex]
\chapter{Our solution}\label{chap:oursolution}

The first attempt was to create wrapper of the \react library into the Dart language. 
This wrapper was successfully created, tested and also used in an independent commercial project. 

The problem occurs in the performance of the wrapper, 
where the bottleneck of the speed was the communication between \react created in the JavaScript and the wrapper in the Dart language. 
This bottleneck can be reduced by some adjustment and some Dart hacks, but it was still a bottleneck. 

That's why we decided to build our own library called \tiles.
Most of the performance benchmarks of the \tiles library later in this work will be compared with the \react wrapper. 

As we told in the previous part of the work, 
we decided to take inspiration from the \facebook \react library,
mainly in the API of the library, which is component based, 
with some differences in architecture.

We don't have to implement some of the additional features of the \react library, because of the nature of the Dart language.
\begin{itemize}
  \item Synthetic events \textit{(Dart unified events)}
  \item Mixins \textit{(Dart support native mixins as a part of a language)}
  \item Props type checking \textit{(Dart is optional-typed language)}
  \item Get default props and initial state \textit{(Dart work with classes which have constructors)}
  \item Changed class name \textit{(Map in dart use string, so string "class" is no more reserved word)}
  \item Test utilities \textit{(Dart has own unittest library and we work with classes and with native events, it is easily tested)}
\end{itemize}

In this next sections of this chapter we will introduce and deeply describe our Dart library \tiles.

\section{Requirements}\label{subsec:our-architecture-requirements}

  When we designed the \tiles library, we take into account requirements derived from the motivation of this work:

  \begin{description}
    \item{\textbf{Rendering in both environments, the browser and the server}} \hfill \\
      One of the main advantages of \tiles library is a possibility to render the same content, 
      with the same code as on the server, so in the browser.
      This resolved into the package structure and several architectural decisions.
      \begin{description}
        \item{\textbf{No template usage}} \hfill \\
          To achieve a possibility to render content in both environments, easy testing and mocking, 
          we decided not to use templates.
      \end{description}
    \item{\textbf{Easy to use API}} \hfill \\
      The solution can be very interesting and powerful, 
      but if it don't offer a reasonably easy to use API, almost no one will use it.
      \begin{description}
        \item{\textbf{\react-like API}} \hfill \\
          Because we use a similar concept as the JavaScript \react library,  
          which is widely used and known, 
          if we offer similar API, more people will quickly get used to it.
      \end{description}
    \item{\textbf{Performance}} \hfill \\
      We want to offer the useful library, and if want someone to use it, 
      we need to offer a good performance in the competition of Dart and JavaScript UI libraries.
  \end{description}

  How we fulfilled these requirements is in detail described in next sections of this chapter.

\section{Architecture - \tiles}\label{sec:our-architecture}

	The \tiles library implements a component driven databinding. 
	It enable declarative definition of the user interface by the \texttt{Component} class.
	The component define the structure of the virtual DOM tree, which is rendered to the real DOM. 
	Updates in the virtual DOM are reflected into the real one with usage of the smallest difference between old and the new structure.
	
	The composition of the virtual DOM is used in the \tiles to create the system of synthetic event bubbling, 
	which offer the programmer a simple way of listening to events in the DOM from a component perspective.

	The architecture of the \tiles library offer the possibility to think about UI as a set of the functional components, 
	instead of HTML elements and a Dart/JavaScript application manipulating with them.

  \subsection{Architectural overview}\label{subsec:our-architecture-overview}
    Our high level idea is based on the \facebook \react library attitude. 
    We created api, whose main class is \texttt{Component}, which represents construct very similar to \react's \texttt{Component}.
    This component is mounted to an element, where it renders itself. 
    This relationship is described on \fullref{img:library-idea}.
    \begin{figure}[h]
    \centering  
      \includegraphics[scale=0.6]{images/uml/tiles/Idea/Idea.png}
      \caption{The Idea}
      \label{img:library-idea}
    \end{figure}

    These components are placed into tree structure, which represents \mbox{\textbf{\nameref{img:library-idea-virtual-dom}}}, 
    which is then translated to the real DOM of client's browser or to the markup rendered by server application.

    \begin{figure}[h]
    \centering  
      \includegraphics[scale=0.6]{images/uml/tiles/Idea/IdeaObjects.png}
      \caption{Virtual DOM}
      \label{img:library-idea-virtual-dom}
    \end{figure}
    There can be the event listeners attached to these components. 
    The \textbf{events} 
    \footnote{We work at Dart, which create browser compatibility for us, so we don't have to create synthetic events like \react.}
    are then bubbled trough a virtual DOM, instead real one.
    By this there can be the listener attached to a custom component, which doesn't have element representing it in a real DOM.

    As we work in Dart language, it is natural to try to reuse the most of code on the both, client and server side.
    The next important part of architecture is \textbf{server-side rendering}.

    It is very important for SEO purposes and smooth user experience. 

  \subsection{Structure}\label{subsec:our-architecture-structure}

    We split our library into 3 partially dependent packages. 
    \begin{description}
      \item[Tiles] \hfill \\
        \textbf{Tiles} package creates the core component's of library, focused to create and maintain virtual DOM and provide API for programmer.
        This package should be included by programmer in the files, where he defines custom components. 
        These components then can be used on both, server and browser sides.
      \item[Tiles Browser] \hfill \\
        This package is used for mounting components to the HTML elements. 
        It maintains relationships between elements and components, 
        simulates events bubbling and keeps real DOM in sync with virtual one.
      \item[Tiles Server] \hfill \\
        \textbf{Tiles Server} package maintains server-side rendering. It offers an API to render component structure to string with markup based on DOM components.
    \end{description}

    Based on the mentioned packages structure, it is quiet obvious what are the dependences between these packages.
    \texttt{\textbf{Tiles}} package is independent, 
    and both of \texttt{\textbf{Tiles Browser}} and \texttt{\textbf{Tiles Server}} are dependent on \texttt{\textbf{Tiles}} package.
    These dependences are shown on \fullref{img:library-packages}.

    \begin{figure}[h]
    \centering  
      \includegraphics[scale=0.6]{images/uml/tiles/Packages.png}
      \caption{Packages}
      \label{img:library-packages}
    \end{figure}

  \subsection{Core}\label{subsec:our-architecture-core}
    There are 4 main classes in the core of the library.
    \begin{description}
      \item[Component] 
        represents closed block of user interface, 
        that should be rendered in application.
      \item[Node]
        is a vertex in a tree of virtual DOM. 
        It contains an instance of \texttt[Component], which represents the type of this \texttt{Node}.
      \item[ComponentDescription]
        is self-explanatory. 
        It is returned from the component to describe it's children. 
        The principles are described later at this document.
      \item[NodeChange]
        represents one change in a virtual DOM, which should be applied into the real DOM. 
        This way, we are able to achieve minimal changing of the real DOM. 
        Types of change are: \texttt{CREATED, UPDATED, MOVED, DELETED}
    \end{description}

    In contrast with \facebook \react \texttt{Component}, our component provides only an API to a programmer.
    This class is the main class for a programmer using our library. 
    He doesn't need to use any other class created by our library. 
    Just some methods.

    We also have an inspiration from \react with idea of \textbf{virtual DOM}. 
    Vertices of an virtual DOM are represented by the class \texttt{Node} instead of the class \texttt{Component} to separate the functionality. 

    Each node contain an instance of \texttt{Component}. 
    The node represents the component in a virtual DOM.
    
    The diagram of relationships is shown on \fullref{img:library-core}. 
    In the next chapters we describe the main classes more in details.
    \begin{figure}[h]
    \centering  
      \includegraphics[scale=0.2]{images/uml/tiles/TilesCore/Core.png}
      \caption{Core of the library}
      \label{img:library-core}
    \end{figure}

    \subsubsection{Component}\label{subsubsec:our-architecture-core-component}
      The Component is the main building brick of application(library). 
      It offers api to the programmer with life-cycles, props etc. 

      Component is a class, which represents functionality of certain part of UI\footnote{UI = User Interface} in an application.
      It is created with some props and children acting as parameters of a function. 
      Main purpose of the component is to create a structure below\footnote{From the virtual DOM tree point of view} it, 
      add the event listeners and update itself based on event listener input.

      The main method of the \texttt{Component} is \texttt{List<ComponentDescription> render()}. 
      By this method component describes it's substructure.
      It will return list of children of this component, represented by instance of the class \texttt{ComponentDescription}. 
      \texttt{Node}, which owns this component (and called it's render method) will manage the rest.
      Basically, it will return the message like \textit{"This is how I should look like"}.

      Second important method is \texttt{void redraw()}, which trigger redraw of the component. 
      This redraw will be executed on the next animation frame. 

      Redraw is powered by \texttt{needUpdate} stream offered by \texttt{Component}, 
      which is automatically created in default constructor of class \texttt{Component}, 
      so it is very important, to call superclass constructor in each custom component class.

      \texttt{ComponentDescriptionFactory registerComponent(ComponentFactory factory)} is additional method helping programmer 
      for easier \texttt{ComponentDescription} creation. 
      Factories are described late in \fullref{sec:our-api}
      
    \subsubsection{DomComponent}\label{subsubsec:our-architecture-core-dom-component}
      \texttt{DomComponent} is a subclass of class \texttt{Component}. 
      This is specialized class, which represents HTML elements in the component structure.

      It has \texttt{props} saved as \texttt{Map}, because HTML element has attributes saved in \texttt{Map}.
      \texttt{render} method returns \texttt{children} member variable and \texttt{svg} and \texttt{pair} flags.

      Specific HTML elements are created based on different \texttt{ComponentFactory} and \texttt{ComponentDescriptionFactory}.
      \texttt{ComponentDescriptionFactory} is used to easily create \texttt{ComponentDescription}s of \texttt{DomComponent} in a custom component render method. 

    \subsubsection{ComponentDescription}\label{subsubsec:our-architecture-core-component-description}
      \texttt{ComponentDescription} is a description of the component. 
      It describes which type of the component should be rendered by using which parameters.

      For this purpose, it needs 4 types of information:
      \begin{itemize}
        \item \textbf{Type of the component} \hfill \\
          To create instance of a component, we need to know, what type (class) of the component it should be. 
          This information is represented by \texttt{ComponentFactory}, 
          which is function with 2 parameters, \texttt{props} and \texttt{children}, 
          which returns instance of a subclass of a \texttt{Component}.
        \item \textbf{Properties} \hfill \\
          Data which should be passed to the factory. 
          This data are used as a properties of the component.
        \item \textbf{Children} \hfill \\
          Children of described component.  
          This is useful mainly when programmer wants to render more complex structure of \texttt{DOMComponent}s.
        \item \textbf{Key} \hfill \\
          Key is an identifier of a child. 
          It is used to recognize reordering of children of the component.
          When components \texttt{render} method returns list of descriptions, 
          keys they contain are recognized and matched with keys stored in virtual DOM.

          If there is a match in key of the child in different position, child is only moved an updated. 
          If there is no match in key, default process follows.
      \end{itemize}

      Description is once created with all the parameters and then these parameters can't be changed. 
      All these parameters is set up by constructor. 

      \texttt{ComponentDescription} has one important method, which is \texttt{Component createComponent()}, 
      which creates \texttt{Component} instance with props and children from the description.

    \subsubsection{Node}\label{subsubsec:our-architecture-core-node}
      \texttt{Node} is the most important and complex class in the library.
      It providers following functionality:
      \begin{itemize}
        \item creates virtual DOM tree, maintains creating and updating of the tree based on results of component's \texttt{render} method, 
        \item listens to component's \texttt{needUpdate} stream and marks self as \textit{"dirty"} when it's component need update,
        \item and handles updating process that is rearranging children of the \texttt{Node}.
      \end{itemize}

      The node is also a vertex of the virtual DOM. 
      It store children as a list of children. 
      To use all possible optimization, node contain a \texttt{ComponentFactory} of the contained component, which is used when the virtual DOM is updated.
      It also contain \texttt{key} which is used to recognize changed position of the same child.

      Node has two important flags: \texttt{isDirty} and \texttt{hasDirtyDescendant}. 
      These flags represent information, whether the node, or its descendants, needs to be redrawn.
      If \texttt{isDirty} is true, the node needs to be updated, 
      because component of this node called \texttt{redraw} method.
      If \texttt{hasDirtyDescendant} is true,  
      there exist a descendant of this node, which wants to be updated.
      When \texttt{hasDirtyDescendant} is true and \texttt{isDirty} is false, 
      the node doesn't have to update itself, it is enough to call update on child nodes.

      Method \texttt{update({List<NodeChange> changes, bool force: false})} is executing the update process. 
      It take 2 named arguments:
      \begin{description}
        \item{\textbf{\texttt{List<NodeChange> changes}}} is used to be filled by changes generated by the update process,
        \item{\textbf{\texttt{bool force: false}}} is a flag signalizing if to update despite that node is not dirty.
      \end{description}

      The \texttt{update} method is used mostly in the browser part of an library, where it offers a possibility to get changes in the virtual DOM, which should be used to update the real one. 

      A methods logic consists of several main steps.
      \begin{enumerate}
        \item check, if update is needed by flags \texttt{isDirty}, \texttt{hasDirtyDescendant} or \texttt{force}, if no, exit,
        \item if component of this node needs to update (\texttt{isDirty == true}) or \texttt{force == true}, update this node with rearrangement of children,
        \item if any changes was generated, add them into the \texttt{changes} list
        \item set this node as not dirty and not have dirty descendant.
      \end{enumerate}

      The algorithm of rearrangement of children by calling \texttt{render} method of this component 
      and adapting node's children to returned descriptions is fully documented in the source code related to this work: 
      \url{https://github.com/cleandart/tiles/blob/master/lib/src/core/node_update_children.dart}

    \subsubsection{NodeChange}\label{subsubsec:our-architecture-core-node-change}
      \texttt{NodeChange} takes place as a record of a change in the virtual DOM. 
      It is used to mirror changes in the virtual DOM with the real DOM. 

      When some node in the virtual DOM is updated by method \texttt{update}, the list of changes is is collected.
      This list is subsequently processed by browser part of a library, which mirrors changes to the real DOM.

      \texttt{NodeChange} class has no methods (except constructor) and acts as a data chunk dedicated just for it's purpose. 
      It contains node, type of change, old and next properties. 

      Type is stored as an instance of \texttt{NodeChangeType} enum and can be one of: \texttt{CREATED}, \texttt{UPDATED}, 
      \texttt{MOVED} and \texttt{DELETED}. 
      When type is \texttt{UPDATED}, old props and new props take effect. 

  \subsection{Life-cycle}\label{subsec:our-architecture-lifecycle}

    Every instance of \texttt{Component} has its own life-cycle. 
    As every object, first it is created. 
    Subsequently, it is mounted and rendered into the virtual DOM, and then in to real DOM.

    When a "higher" node wants to update the node of the component, 
    the component will first receive props, 
    subsequently is asked if want to update, 
    if yes, it is asked for actual description of its children. 

    Sometimes the component wants to update itself (e.g. because event occurs).
    It calls redraw, then, it will be asked if really should update, and if yes, it is rendered and updated.

    At the end of component's life in the real DOM, 
    component should be notified about that.

    The whole life-cycle is shown on the \fullref{img:our-life-cycle}.

    \begin{figure}[h]
    \centering  
      \includegraphics[scale=0.5]{images/uml/tiles/TilesCore/LifeCycle.png}
      \caption{Life cycle of a \texttt{Component}}
      \label{img:our-life-cycle}
    \end{figure}

    \subsubsection{Create}\label{subsec:our-architecture-lifecycle-create}

      The create part of life-cycle is implemented by constructor of \texttt{Component}. 
      It will receive props and optionally children as arguments and 
      it should prepare the whole state of object to live.

      An trivial example of constructor of \texttt{Component} is displayed below.
      % flatex input: [content/examples/my_component_constructor.tex]
      \begin{minted}{dart}
class MyComponent extends Component {
  MyComponent(props, [children]): super(props, children) {}
}
      \end{minted}

% flatex input end: [content/examples/my_component_constructor.tex]


      Constructor of an example only calls constructor of superclass \texttt{Component}.

      Example of more complex constructor should be e.g. the \texttt{Todo} component example:
      % flatex input: [content/examples/my_todo_component_constructor.tex]
      \begin{minted}{dart}
class MyTodoComponent extends Component {
  Todo todo;
  MySearchComponent(props, [children]): super(props, children) {
    if (props != null && props.todo is Todo) {
      this.todo = props.todo
    } else {
      this.todo = new Todo();
    }
  }

  // ...

}
      \end{minted}

% flatex input end: [content/examples/my_todo_component_constructor.tex]


    \subsubsection{Did mount}\label{subsec:our-architecture-lifecycle-didmount}

      When the component is mounted to the real DOM, user of the library should be notified about this event. 
      It is done by triggering the \textbf{Did mount} life-cycle event implemented the method \texttt{didMount}. 

      This is the place in time, were the component \textit{start to live its life} with the connection to the real DOM.
      This is the correct place to initialize for example the timers, stream listeners etc. 

      Our \texttt{MyTodoComponent} example compoent should listen for a change of the todo on the server, 
      and if it was changed, we can redraw the component.
      % flatex input: [content/examples/my_todo_component_didmount.tex]
      \begin{minted}{dart}
class MyTodoComponent extends Component {
  Todo todo;
  StreamSubscription subscription;

  // ...

  didMount() {
    this.subscription = this.todo.changedOnServer.listen((change) {
      this.redraw();
    });
  }

  // ...

}
      \end{minted}

% flatex input end: [content/examples/my_todo_component_didmount.tex]


    \subsubsection{Will receive props}\label{subsec:our-architecture-lifecycle-willreceiveprops}

      When the component is updated by "higher" ancestor, it will receive new props.
      A user of the library can need the possibility to compare these props with the old one
      and perform needed changes. 

      This is the correct place for subsequent life-cycle event \textit{Will receive props}
      implemented by the method \texttt{willReceiveProps}. 

      The example of \texttt{willReceiveProps} in \texttt{MyTodoComponent} 
      should compare the \texttt{todo} of the old and the new props and if they are not equal, 
      update change listener.
      % flatex input: [content/examples/my_todo_component_willreceiveprops.tex]
      \begin{minted}{dart}
class MyTodoComponent extends Component {
  Todo todo;
  StreamSubscription subscription;

  // ...

  willReceiveProps(dynamic newProps) {
    if (this.todo != newProps.todo) {
      this.subscription.cancel();
      this.subscription = newProps.todo.changedOnServer.listen((change) {
        this.redraw();
      });
    }
  }

  // ...
}
      \end{minted}


% flatex input end: [content/examples/my_todo_component_willreceiveprops.tex]

      
    \subsubsection{Should update}\label{subsec:our-architecture-lifecycle-shouldupdate}

      An optimization of the performance of an application can be done by rejecting "redraw" of the component.
      To reject this "redraw", component should be asked, if the redraw is needed. 

      This is implemented by the \texttt{shouldUpdate} method, which returns true if the component want to be redrawn. 

      By default \texttt{shouldUpdate} returns true, 
      what resolves to always updating of custom component, which doesn't implement this method.

      In a basic scenario this method recognize, 
      if the component will be rendered differently with the new props.  
      If not, it return false, else it return true.

      Example in \texttt{MyTodoComponent}:
      % flatex input: [content/examples/my_todo_component_shouldupdate.tex]
      \begin{minted}{dart}
class MyTodoComponent extends Component {
  Todo todo;
  StreamSubscription subscription;

  // ...

  shouldUpdate (newProps, oldProps) {
    if (newProps.todo == oldProps.todo) {
      return false;
    }
    return true;
  }

  // ...

}
      \end{minted}


% flatex input end: [content/examples/my_todo_component_shouldupdate.tex]


    \subsubsection{Render}\label{subsec:our-architecture-lifecycle-render}

      Render is the main part of the \texttt{Component}. 

      It is implemented by the method \texttt{render}.
      It should return array of component descriptions which should be considered as 
      \textit{"this is how this component should look like"}.

      For example, in our \texttt{MyTodoComponent} render will return \texttt{<div>} which contains title and description of \texttt{todo}.
      % flatex input: [content/examples/my_todo_component_render.tex]
      \begin{minted}{dart}
class MyTodoComponent extends Component {
  Todo todo;
  StreamSubscription subscription;

  // ...

  render () {
    return div ({"class": "todo"}, [
      h2 ({}, todo.title),
      p ({}, todo.description)
    ]);
  }

  // ...

}
      \end{minted}


% flatex input end: [content/examples/my_todo_component_render.tex]


    \subsubsection{Did update}\label{subsec:our-architecture-lifecycle-didupdate}

      When the life-cycle method \texttt{didUpdate}, 
      by which is implemented the  life-cycle \textit{"Did update"} event, 
      is triggered, the component is notified, 
      that it is mounted and there exist HTML elements in the DOM for each DomComponent descendant.

    \subsubsection{Will unmount}\label{subsec:our-architecture-lifecycle-willunmount}

      The \textit{"Will unmount"} event is implemented by the method \texttt{willUnmount}.
      It is called right before the component is removed from the virtual DOM and therefore from the real one. 
      This is the correct place to stop all timers and listeners.

      % flatex input: [content/examples/my_todo_component_willunmount.tex]
      \begin{minted}{dart}
class MyTodoComponent extends Component {
  Todo todo;
  StreamSubscription subscription;

  // ...

  willUnmount () {
    subscription.cancel();
  }

  // ...

}
      \end{minted}

% flatex input end: [content/examples/my_todo_component_willunmount.tex]


  \subsection{Rendering}\label{subsec:our-architecture-rendering}

    The main target of the \tiles library is the rendering of a content. 
    By the \texttt{Component} and the \texttt{Node} we can create the virtual DOM tree, 
    which should be reflected into the real one, or into the HTML markup string. 

    As we described earlier in the \fullref{subsec:our-architecture-structure}, 
    these rendering types are separated to separate packages from the core package and are independent from each other. 

    \subsubsection{Server side}\label{subsec:our-architecture-rendering-server}
      On the server, we don't have DOM elements available, 
      therefore we want to render our virtual DOM structure into the string
      representing HTML markup of the virtual and also the real DOM.

      The markup string constructed on the server can be reused in the browser to smooth user experience. 

      Our target, render virtual DOM into the markup string is not so complicated.
      From the programmers point of the view, 
      he will use a \texttt{ComponentDescription} to describe the component to be rendered.
      Our server-side package of the library will get the description and construct the node tree structure, 
      with node containing described component in the root. 
      The node tree represents the created virtual DOM. 

      From the virtual DOM we can render corresponding markup string by the depth-first search of it's tree. 
      The markup string is created recursively for every subtree starting from a node by this algorithm:

      \begin{algorithm}[H]
        \KwData{The node in virtual DOM tree}
        \KwResult{A string with the markup of the virtual DOM subtree with a root in the node}
        \eIf{the component of the node is DomComponent}{
          \eIf{the component is not pair}{
            write the markup into the result with attributes from the props and a tag name from the component\;
          }{
            write the open markup into the result with attributes from the props and a tag name from the component\;
            write the markup for all children recursively into the result\;
            write the close markup into the result\;
          }
        }{
          write the markup for all children recursively into the result\;
        }
        % }
        \caption{Write node into the markup string.}
      \end{algorithm}

    \subsubsection{In browser}\label{subsec:our-architecture-rendering-browser}

      The rendering in the browser is more difficult than the rendering into the markup string. 
      It is possible to use the same \textit{render to the markup string} method and add this markup to the DOM. 
      This will create DOM structure representing the virtual DOM, 
      but doesn't create connections between nodes in the virtual DOM and HTML elements in the real one. 
      This connections is necessary when the virtual DOM is updated.
      Therefore the browser side rendering will work in a different way. 

      \paragraph{Initial mount}\label{par:our-architecture-rendering-browser-init} 

        The first thing, the user of the library need to do, is to mount the component into the HTML element(\textit{mount root}). 
        For this purpose the \texttt{ComponentDescription} describing the component is used.
        The mount process consist from creation of the component from the description, 
        placing it into the node and first render of the virtual DOM with the root in the node. 

        The virtual DOM is subsequently reflected into the real DOM created under the \textit{mount root}.

        Suppose that the \textit{mount root} have no children 
        (the case of the \textit{mount root} with children is described in the \fullref{subsec:our-architecture-injecting},%\\
        the node mount process is described in the \fullref{algorithm:mount-node}:
        \begin{algorithm}[H]
          \KwData{The node in the virtual DOM and a HTML element, to mount the node to (mountRoot)}
          \KwResult{Mounted the node into the element}
          \eIf{the component of the node is DomComponent}{
            create an element representing the component\;
            add the created element to the mountRoot\;
            save relations between the created element and the node;

            \If{the component is a pair component}{
              \For{a child in children of the node}{
                mount the child node into the newly created element as a mountRoot\;
              }
            }
          }{
            \eIf{the component is a text component}{
              create a HTML text node with the text from the component\;
              add the text node to the mountRoot\;
              save relations between the created text node and the node;
            }{
              \For{a child in children of the node}{
                mount the child node into the mountRoot\;
              }
            }
          }
          % }
          \caption{Mount the component into the real DOM.}
          \label{algorithm:mount-node}
        \end{algorithm}

        The algorithm is recursive and skips custom components. 
        It creates relationships between created elements and nodes. 
        These relations is used to easy reflect the change in the virtual DOM into the real one 
        when the virtual DOM is changed.

        The DOM structure created by the \fullref{algorithm:mount-node} reflects the virtual DOM structure 
        with removed nodes with the custom component.

      \paragraph{Update}\label{par:our-architecture-rendering-browser-update} 
        
        As time goes, the situation, that the virtual DOM need to be updated, occurs. 
        The need of the update is realized by the method \texttt{redraw} of the component, 
        which cause marking the node as dirty.
        The framework recognize this situation and perform update of this node and the structure under it.
        The update produces a list of changes in the virtual DOM needed to be applied to the real DOM.

        Updates are processed by their type. 
        For every type the information about the HTML element representing a node is needed. 
        This information is stored in the relation \mbox{\texttt{Node $\rightarrow$ Element}}, 
        created in the mount of the node. 
        It is stored by \texttt{Map<Node, Element>}.

        What happened when we want to apply the node change into the real DOM structure differs by the type of the change:
        \begin{description}
          \item[CREATED] \hfill \\
            A newly created node is mounted into the DOM. 
            The adjusted \fullref{algorithm:mount-node} is used for this purpose.
          \item[UPDATED] \hfill \\
            If the node with a DOM component in it was updated, the element, related with it, is updated with the props of the node.
            If the node with a custom component was updated, no action is needed, because this node is not directly reflected to the DOM.
          \item[MOVED] \hfill \\
            The DOM representation of the moved node is moved to the new position.
          \item[DELETED] \hfill \\
            The DOM representation of the removed node is removed from the real DOM. 
        \end{description}

  \subsection{Events}\label{subsec:our-architecture-events}

    In the library which create own virtual DOM structure, it is natural to think about event in this virtual DOM. 
    As in the real one, it is natural to create event bubbling trough virtual DOM. 
    This synthetic event bubbling represents real event bubbling of the event 
    trough the real DOM elements associated with nodes in the virtual DOM.

    Synthetic event bubbling is useful for the user of the library to catch events in the DOM and react on them by update of the state 
    and triggering of the redraw of the component, if needed. 

    The \tiles library listen to DOM events on \textit{mount root} of the virtual DOM.
    When an event occur, the \tiles catch it on it's bubbling route up to he \texttt{body} element, 
    recognize, which node represents the target element of the event and start synthetic bubbling of the event from that node. 

    To store the relation \mbox{\texttt{Target Element $\rightarrow$ Node}} is used \texttt{Map<Element, Node>}.

    The \tiles library listen to every type of the event only once in the \textit{mount root} 
    of the virtual DOM for each event listened by components in it.

    It is important from the performance point of view, because we can add only one event listener for each event type in the whole virtual DOM.
    We will discuss this later in \fullref{chap:performance}.

    Synthetic event bubbling trough the virtual DOM enable to listen to the bubbled event even on custom component.

    \subsubsection{Synthetic bubbling}\label{subsec:our-architecture-events-bubbling}

      When the component is mounting, we store the relation between the HTML element, and the node, which contains this component. 
      Then we check a presence of the event listeners in the descriptions of the component. 
      We add event listener for every type of the event present in the virtual DOM to the \textit{mount root} HTML element. 
      The event listener in the \textit{mount root} is only one for each event type in the virtual DOM.

      When the event type listened by \textit{mount root} occurs in the DOM subtree representing our virtual DOM, 
      it will bubble up to the \textit{mount root}, where is caught by the event listener created by the \tiles library. 
      The \tiles library will recognize the target HTML element of the caught event 
      and simulate event bubbling from the node representing the target element in the virtual DOM.

      The synthetic bubbling process starts from the target node and check event listeners associated with it.
      If a listener for a type of the current event is presented, it is called with the event and the component as arguments. 

      The component is passed because the listener is not created by this component itself, 
      but by the component above it, which should be informed, on which component the listener was called.

      There is no official way of getting an information if the \texttt{stopPropagation} was called on the event.
      Therefore the listener, which want to stop a propagation, should return the \texttt{false} boolean.
      If the listener returns something else, bubbling continues with the parent of the current node.
      When the root node is achieved, the bubbling stops.

  \subsection{Injecting}\label{subsec:our-architecture-injecting}

    The \tiles library implement a built in server side rendering.

    It is important as for the SEO, so for the user experience. 
    User don't have to wait until the source code of the application is loaded into the browser and see the content of the page. 
    If our browser package replaces the structure generated on server with a new DOM structure 
    generated from virtual DOM, the part of the page is re-drawn from scratch,
    what cause disappearance of the content for the small amount of the time.

    To prevent this behavior, we created an injecting system, which will "inject" the existing DOM structure 
    and rebuild it to represent the virtual DOM.

    When the \texttt{ComponentDescription} is mounting, 
    basic implementation can erase whole content of element to which is the description mounted to.
    Instead, we will reuse the existing structure by iterating trough the virtual DOM and reusing every element, which match the virtual DOM. 
    The \tiles library iterate trough the virtual DOM and the real DOM simultaneously and compare them. 

    When the iteration trough the virtual DOM process a node with a DOM component and the HTML element,
    their tag names are compared.
    If tag names match, the HTML element is associated with the node, 
    element's attributes are adapted to the props of the node's component 
    and is used to mount/inject children of the node recursively.

    When the tag name don't match, 
    a new HTML element is inserted before the currently iterated element and paired with the node. 
    Subsequent DOM component at the same layer of the DOM component tree
    \footnote{
      By the DOM component tree we mean a derived tree from the virtual DOM tree, 
      which can be constructed by removing nodes with a custom component
      and connecting nodes with DOM components to the closest ancestor in the virtual DOM tree with a DOM component in it.
    }
    can reuse this not matching element.

    When injecting of children finishes for a node with a DOM component in it, 
    not used(injected) child HTML elements of the associated element are removed.

    This implementation of injecting enable a full reuse of the DOM structure created on the server.
    It also create relations important for future updates of the virtual DOM. 

    Moreover, if there exist a similar structure, not generated from same components on the server, 
    our library will reuse as much as possible of it.

\section{API}\label{sec:our-api}

  One product of this work is an open source UI library called \tiles.
  To use this library, user need to know the application programming interface(API) of it. 

  In the \tiles library, the user interface consist of components, represented by the class \texttt{Component}. 
  The \texttt{Component} class is used to create a functional logic of a part of the UI. 

  To use it, mount it and add it to the other component as a child is necessary to create a \texttt{ComponentFactory} 
  and a \texttt{ComponentDescriptionFactory}, which creates a \texttt{ComponentDescription}.

  The \textbf{\texttt{ComponentFactory}} is a function with two positional optional arguments: \texttt{props} and \texttt{children}.
  It return an instance of the \texttt{Component}. The \texttt{ComponentFactory} is used as a comparable type of the component.
  The easiest and most common way to create \texttt{ComponentFactory} is shown in next example:
  % flatex input: [content/examples/component_factory.tex]
\begin{minted}{dart}
class MyComponent extends Component {
  MyComponent(props, [children]): super(props, children) {}
}

ComponentFactory factory = 
  ([dynamic props, dynamic children]) => new MyComponent(props, children);
\end{minted}

% flatex input end: [content/examples/component_factory.tex]

% }

  \textbf{\texttt{ComponentDescription}} is a description of a component. 
  It is used as argument to mounting functions and as a result of the \texttt{render} method.
  Most of the time is used indirectly by usage of \texttt{ComponentDescriptionFactory}.

  \textbf{\texttt{ComponentDescriptionFactory}} is a function, which create a \texttt{ComponentDescription}.
  It has 7 optional named arguments: 
  \begin{description}
    \item[props] passed to the component factory,
    \item[children] also passed to the component factory, 
    \item[key] used to optimize reorder of children and
    \item[listeners] used to attach event listeners to the component created by the description.
  \end{description}

  The \texttt{ComponentDescriptionFactory} can be created by programmer himself, 
  but the most often way is to create it by the \texttt{registerComponent} function, which accept one argument of \texttt{ComponentFacotry} type.
  This method creates \texttt{ComponentDescriptionFactory}, which returns a description with passed \texttt{ComponentFacotry}.
  A usage is shown on the following example:
  % flatex input: [content/examples/register_component.tex]
\begin{minted}{dart}
ComponentDescriptionFactory myComponent = registerComponent(factory);
\end{minted}

% flatex input end: [content/examples/register_component.tex]

% }

  The \texttt{ComponentDescriptionFactory} is used to create a component description or directly in the mounting into HTML element:
  % flatex input: [content/examples/mount_component_description_factory.tex]
\begin{minted}{dart}
Element mountRoot = querySelector("#container");

var props = {};
var children = [];

ComponentFacotry myComponentDescription = myComponent(props, childre);

// Or directly in mountComponent

mountComponent(myComponent(props, children), mountRoot);
\end{minted}

% flatex input end: [content/examples/mount_component_description_factory.tex]

% }

  \subsection{Component}\label{subsec:our-api-component}

    The \texttt{Component} class is the main class of the API of the \tiles library.
    Every custom component should extend or implement it. 

    The \texttt{Component} class contains a constructor, life-cycle methods, 
    the \texttt{render} and a \texttt{redraw} method, props, children and offer a \texttt{needUpdate} stream.
    The default \texttt{Component} implementation is following: 
    % flatex input: [content/examples/component.tex]
\begin{minted}{dart}
class Component {
  dynamic props;
  List<ComponentDescription> children;
  final StreamController _needUpdateController;
  Stream<bool> get needUpdate => _needUpdateController.stream; 

  /**
   * Life cycle
   */
  Component(this.props, [this.children]): 
    this._needUpdateController = new StreamController<bool>() {}
  didMount() {}
  willReceiveProps(dynamic newProps) {}
  shouldUpdate(dynamic newProps, dynamic oldProps) => true;
  List<ComponentDescription> render() {}
  didUpdate() {}
  willUnmount() {}
  
  redraw([bool now = false]) {
    _needUpdateController.add(now);
  }
}
\end{minted}

% flatex input end: [content/examples/component.tex]

% }

    The easier way to create own component is by extending the \texttt{Component} class:
    % flatex input: [content/examples/my_component.tex]
\begin{minted}{dart}
class MyComponent extends Component {
  MyComponent(props, [children]): super(props, children) {}

  render() {
    return div(null, [
             span(null, "This is simple component")
           ]);
  }
}
\end{minted}

% flatex input end: [content/examples/my_component.tex]

% }

    The \texttt{MyComponent} only renders \texttt{ComponentDescription} of a div DOM component which contain children with a span description.

  \subsection{DOM component API}\label{subsec:our-api-dom}

    The \tiles library contain prepared \texttt{ComponentDescriptionFactory} functions for the each of standard HTML elements, 
    with standard arguments. 
    They creates new \texttt{ComponentDescription}, describing \texttt{DomComponent} with passed props and children.

    Usage of the factories are shown on the previous example.

  \subsection{Browser specific API}\label{subsec:our-api-browser}
    
    There are 3 main additions when the virtual DOM is rendered in the browser.
    \begin{itemize}
      \item Mounting
      \item References
      \item Event listeners
    \end{itemize}

    The \textbf{mounting} is implemented by the \texttt{mountComponent} method with 2 arguments, 
    \texttt{ComponentDescription description} and \texttt{Element mountRoot}.
    It will mount the described component into the \texttt{mountRoot} element.
    % flatex input: [content/examples/mount_component_description_factory_only.tex]
\begin{minted}{dart}
mountComponent(myComponent(props, children), mountRoot);
\end{minted}

% flatex input end: [content/examples/mount_component_description_factory_only.tex]

% }

    \textbf{References} are part of the props. 
    If component have in \texttt{props["ref"]} instance of an internal class \texttt{{\textunderscore}Ref}, 
    which is only a function returning \texttt{void} with one \texttt{Component} argument, 
    then when this component is created and mounted, 
    this reference method is called with it. 

    It is useful, when a custom component want to have a reference to the element associated to one of it's descendant. 
    Example of the usage is in the following example:
    % flatex input: [content/examples/ref.tex]
\begin{minted}{dart}
class MyComponent extends Component {
  /* ... */
  Element input;

  render() {
    return input({"ref": (component){
      this.input = getElementForComponent(component);
    }});
  }
}
\end{minted}

% flatex input end: [content/examples/ref.tex]

% }

    \textbf{Event listeners} are represented in a form of a Map, stored at the node. 
    They should be instances of a \texttt{EventListener} class, which is a function with 2 arguments, an event and a component, and returns boolean.
    The map containing event listeners is passed to the \texttt{ComponentDescriptionFactory} as a named parameter \texttt{listeners}.
    % flatex input: [content/examples/event_listener.tex]

\begin{minted}{dart}
input(listeners: {"onClick": (event, component){
  print("Input clicked.");
}});
\end{minted}

% flatex input end: [content/examples/event_listener.tex]

% }

  \subsection{Server specific API}\label{subsec:our-api-server}

    There is only one thing we want to do on the server. Create markup for a \texttt{ComponentDescription}.
    We created a method \texttt{mountComponentToString}, 
    which require one \texttt{ComponentDescription} argument and returns the markup string representing the virtual DOM created from the description of a component.
    % flatex input: [content/examples/mount_component_to_string.tex]
\begin{minted}{dart}
String markup = mountComponentToString(
  span(props: {"class": "my-span"}, "Text in the span")
);
markup == '<span class="my-span">Text in the span</span>' // is true
\end{minted}

% flatex input end: [content/examples/mount_component_to_string.tex]

% }

% flatex input end: [content/chapters/our.tex]

% flatex input: [content/chapters/performance.tex]
\chapter{Performance}\label{chap:performance}

% flatex input end: [content/chapters/performance.tex]

% flatex input: [content/chapters/benchmarks.tex]
\chapter{Benchmarks}\label{chap:benchmarks}


% flatex input end: [content/chapters/benchmarks.tex]

\pagestyle{plain}
% flatex input: [content/chapters/conclusion.tex]
%------------------------------------------------------------------------------------------------------
%  tak isto ako ufod, toto je zaver. V pripade anglictiny premenujte kapitolu...
%------------------------------------------------------------------------------------------------------
\chapter*{Conclusion}\label{chap:conclusion}
\addcontentsline{toc}{chapter}{Conclusion}

Here will be conclusion of wholw thesis


% flatex input end: [content/chapters/conclusion.tex]


% flatex input end: [./content/chapters/chapters.tex]

%------------------------------------------------------------------------------------------------------

\backmatter

%------------------------------------------------------------------------------------
% vyprodukujeme bibliografiu zo suboru literatura.bib
%------------------------------------------------------------------------------------
\nocite{*}
%*flatex input: [main.bbl]
\begin{thebibliography}{INC}

\bibitem[col]{react}
A~Facebook \&~Instagram collaboration.
\newblock {\em {\sl React - A JAVASCRIPT LIBRARY FOR BUILDING USER
  INTERFACES}}.
\newblock \hfill\break{\tt \url{http://facebook.github.io/react/}}.

\bibitem[Goo]{angular}
Google.
\newblock {\em {\sl AngularJS - HTML enhanced for web apps!}}
\newblock \hfill\break{\tt \url{http://angularjs.org/}}.

\bibitem[INC]{ember}
TILDE INC.
\newblock {\em {\sl ember - A framework for creating ambitious web
  applications}}.
\newblock \hfill\break{\tt \url{http://emberjs.com/}}.

\bibitem[Kea]{runloop}
Tyler Keating.
\newblock {\em {\sl The Run Loop}}.
\newblock \hfill\break{\tt
  \url{http://blog.sproutcore.com/the-run-loop-part-1/} and
  \url{http://blog.sproutcore.com/the-run-loop-part-2/}}.

\end{thebibliography}

% flatex input end: [main.bbl]
%FLATEX-REM:\bibliographystyle{alpha}
%FLATEX-REM:\bibliography{literatura.bib}

\addcontentsline{toc}{chapter}{Bibliography}


\end{document}
 

% flatex input end: [main.tex]
